\chapter{有心运动}

\section{一般有心力作用下的运动}

\begin{problem}{4.1.1}
一质量为 $m$ 的质点沿着由 $x = x_0\cos \omega_1t, y = y_0\sin \omega_2t$ 给定的轨道运动.

(1) 求作用力的 $x$ 分量和 $y$ 分量, 问在什么情况下这个力是有心力;

(2) 求势能 $V(x, y)$ ;

(3) 求动能, 证明质点的机械能守恒.
\end{problem}
\clearpage


\begin{problem}{4.1.2}
一个质量为 $m$ 的粒子在大小为 $\frac{2m}{r^3}$ 的有心引力作用下运动. $t = 0$ 时, $r = 2$ , 速度的径向和横向分量分别为 $\frac{\sqrt{3}}{2}$ 和 1. 证明 $\ddot{r} = \frac{2}{r^3}$ , 并求 $r(t)$ .
\end{problem}
\clearpage


\begin{problem}{4.1.3}
一质量为 $m$ 的质点以初速率 $v$ 从无穷远沿一直线运动, 力心 $P$ 离此直线的最短距离为 $a$ , 有心引力使质点在轨道 $r = c \coth \varphi$ 中运动, 求 $F(r)$ 和 $\varphi(t)$ .
\end{problem}
\clearpage


\begin{problem}{4.1.4}
求使一个粒子在 $r = a(1 + \cos \varphi)$ 的轨道上运动的有心力.
\end{problem}
\clearpage


\begin{problem}{4.1.5}
一个质量为 $m$ 的质点在有心力 $F = -\frac{k}{r^3}$ ( $k$ 为正值常量)作用下运动,选择总能 $E$ 和角动量 $J$ 的什么值时,它的轨道有 $r = ae^{b\varphi}(a,b$ 为常量)形式.
\end{problem}
\clearpage


\begin{problem}{4.1.6}
求质量为 $m$ 的质点在具有势能 $V = \frac{\alpha}{r} - \frac{\beta}{r^2} (\alpha, \beta$ 均为正值常量)的有心力场中运动的轨道，并说明在什么运动条件下，质点将通过力心、趋于力心和飞往无限远，有无可能做周期运动？
\end{problem}
\clearpage


\begin{problem}{4.1.7}
一个质量为 $m$ 的质点在有心力 $F = -\frac{mk^2a^4}{r^3} (k, a$ 均为常量)作用下运动。开始位于 $r = a, \varphi = \pi$ ，速度的径向分量和横向分量分别为 $-\frac{ka}{\pi}$ 及 $ka$ ，求：

(1) 轨道方程 $r(\varphi)$ ;

(2) 运动学方程 $r(t)$ 和 $\varphi(t)$ .
\end{problem}
\clearpage


\begin{problem}{4.1.8}
一个单位质量的物体在一有心力场中运动, 它的轨道为 $r = ae^{-b\varphi}, a, b$ 均为常量. $t = 0$ 时, 位于离力心 $a$ 处, 径向速度为 $k$ . 求有心力的势能.
\end{problem}
\clearpage


\begin{problem}{4.1.9}
质量为 $1\mathrm{g}$ 的质点在有心力的吸引下运动, 此力与质点到力心的距离的立方成正比, 在距离等于 $1\mathrm{cm}$ 时引力等于 $1 \times 10^{-5}$ $\mathbf{N}, t = 0$ 时, 质点与力心的距离为 $r_0 = 2\mathrm{cm}$ , 速度 $v_0 = 0.5\mathrm{cm} \cdot \mathrm{s}^{-1}$ , 其方向与由力心到此质点之直线成 $45^{\circ}$ 角. 求此质点的运动轨道及运动学方程.
\end{problem}
\clearpage


\begin{problem}{4.1.10}
质量为 $m$ 的质点受到的有心力为 $F = -m\left(\frac{\mu^2}{r^2} + \frac{\nu^2}{r^3}\right)$ , $\mu, \nu$ 都是常数, 证明: 当 $h^2 > \nu^2 (h$ 为面积速度的两倍)时, 其轨道方程有

\[
r = \frac {a}{1 + e \cos k \varphi}
\]

的形式，并给出 $a, e, k$ 与 $h, r_{\min}$ 等的关系.
\end{problem}
\clearpage


\begin{problem}{4.1.11}
质量为 $m$ 的质点 $A$ 在光滑的水平桌面上运动, 此质点系在一根可不计质量的、不可伸长的绳子上, 绳子穿过桌面上的一个光滑小孔 $O$ 后于另一端系一个相同质量的质点 $B$ . 若质点 $A$ 在桌面上离 $O$ 点 $a$ 处, 沿垂直于 $OA$ 的方向以速度 $\left(\frac{9}{2} ag\right)^{1/2}$ 射出. 证明此质点在以后的运动中离 $O$ 点的距离必在 $a$ 及 $3a$ 之间.
\end{problem}
\clearpage


\begin{problem}{4.1.12}
导出求有心力的公式

\[
F = \frac {m h ^ {2}}{2} \cdot \frac {\mathrm{d} p ^ {- 2}}{\mathrm{d} r}
\]

式中 m 为质点的质量, r 为质点到力心的距离, h 为角动量与质量之比, p 为力心到轨道切线的垂直距离.
\end{problem}
\clearpage


\begin{problem}{4.1.13}
一个在有心力作用下的粒子的运动轨道为 $r\varphi =$ 常量，求其势能 $V(r)$
\end{problem}
\clearpage


\begin{problem}{4.1.14}
稠密气体团质量为 $m$ ，半径极小，可以忽略，像粒子一样，穿过球状的、半径为 $R$ 的低密度的均匀气体云，气体云总质量为 $M(M \gg m)$ ，忽略所有非引力作用。

（1）用“粒子”与云中心的距离 r 表示“粒子”受的力 $F(r)$ 及其势能 $V(r)$ ， $(0 \leqslant r \leqslant \infty)$ ，画出 $V(r)$ 的示意图；

(2) 若“粒子”具有角动量 $J = m\left(\frac{GMR}{32}\right)^{1/2}$ , 总机械能 $E = -\frac{5GMm}{4R}$ , 求轨道的转折点 $r_{\max}, r_{\min}$ , “粒子”总是在云里? 云外? 还是有时在云里, 有时在云外?

(3) “粒子”具有(2)问所给的 $J$ 和 $E$ , 求轨道方程 $r(\varphi)$ , 并画出示意图, 说明有几个转折点.
\end{problem}
\clearpage


\begin{problem}{4.1.15}
一个质量为 $m$ 的质点束缚在线性势 $V = kr(k$ 为常量)之中.

（1）当能量、角动量为多大时，轨道是一个以原点为圆心、半径为 $R$ 的圆？

(2) 圆运动的角速度多大?

(3) 如果质点稍微偏离圆轨道, 小振动的角频率多大?
\begin{center}
\begin{tikzpicture}[>=Stealth, thick]
  % Coordinate axes
  \draw[->] (-2.2, 0) -- (2.4, 0) node[below] {基轴};
  \draw[->] (0, -1.8) -- (0, 1.9);

  % Tilted ellipse and its principal axes
  \begin{scope}[rotate=42]
    \draw (0, 0) ellipse (2.05 and 0.88);
    \draw[dashed] (-2.45, 0) -- (2.45, 0);
    \draw[dashed] (0, -1.45) -- (0, 1.45);
  \end{scope}
\end{tikzpicture}
\end{center}
\end{problem}
\clearpage


\begin{problem}{4.1.16}
质量为 $m$ 的行星围绕质量为 $M$ 的恒星做半径为 $R$ 的准圆周运动，除了万有引力外，此行星还受到一个与行星到恒星的距离成正比的斥力 $f = Ar$ ，试计算行星近心点的运动角速度。
\end{problem}
\clearpage


\begin{problem}{4.1.17}
考虑质量为 $m$ 的行星环绕质量为 $M$ 的太阳的运动轨道，假定有均匀分布密度为 $\rho$ 的尘埃遍布在太阳和行星周围的空间。

(1) 证明尘埃的影响是增加一个有心引力 $F' = -mkr$ ，其中 $k = \frac{4\pi}{3}\rho G, G$ 为万有引力恒量（忽略任何与尘埃碰撞的阻力）；

(2) 若行星在角动量为 $J$ 的圆轨道上运动, 写出轨道半径 $r_0$ 所满足的方程 (不必解).

(3) 假定 $F'$ 与太阳引力相比是小的, 行星的轨道稍微偏离(2)问中所述的圆轨道, 由径向和方位运动的角频率证明轨道是一进动的椭圆, 求出进动角频率 $\omega_{p}$ 与 $r_{0}$ 、 $\rho$ 、G 和 M 的关系;

(4) 椭圆轴运动与轨道运转方向同向还是反向?
\end{problem}
\clearpage


\begin{problem}{4.1.18}
(1) 质量为 $m$ 的行星绕质量为 $M$ 的恒星运动。由于恒星的大气层的阻力，行星受到一个小的阻力 $\bm{F}' = -\alpha \bm{v}$ ，其中 $\alpha$ 为正值常量。设在 $t = 0$ 时行星的轨道是半径为 $r_0$ 的圆，求行星的矢径 $r$ 与时间 $t$ 的函数关系；

(2) 忽略上述阻力,但加上另一小的有心力,行星的势能改写为

\[
V (r) = - \frac {G M m}{r} + \frac {\varepsilon}{r ^ {2}}
\]

点间的幅角 $\varphi$ .

其中 $\varepsilon$ 为一小常量. 可以设想轨道是近似圆形的, 求图4.7中两远心
\begin{center}
\begin{tikzpicture}[>=Stealth, thick]
  % Outer dashed circle representing apocenter envelope
  \draw[dashed] (0, 0) circle (2.1);

  % Precessing rosette orbit: r(t) = p / (1 - e*cos(t)), theta(t) = 65 - (13/12)*t
  % Apocenters at t = 0, 360, 720, 1080 -> theta = 65, 35, 5, -25 deg
  \draw[smooth, domain=-180:1360, samples=400]
    plot ({(0.336 / (1 - 0.84*cos(\x))) * cos(65 - 1.0833*\x)},
          {(0.336 / (1 - 0.84*cos(\x))) * sin(65 - 1.0833*\x)});

  % Angular separation phi between two consecutive apocenters (35 deg and 5 deg)
  \draw[thin] (35:2.1) -- (35:2.55);
  \draw[thin] (5:2.1) -- (5:2.55);
  \draw[<->, thin] (35:2.42) arc (35:5:2.42) node[midway, right=2pt] {$\varphi$};
\end{tikzpicture}
\end{center}
\end{problem}
\clearpage


\begin{problem}{4.1.19}
一粒子在有心引力作用下做半径为 $R$ 的圆周运动，证明当 $f(R) > -\frac{R}{3} \frac{\mathrm{d} f(r)}{\mathrm{d} r} \Bigg|_{r = R}$ 时轨道是稳定的，这里 $f(r)$ 是作用力的大小，它是粒子到力心的距离 $r$ 的函数.
\end{problem}
\clearpage


\begin{problem}{4.1.20}
一质量为 $m$ 的质点在有心引力 $f(r)$ 作用下运动.

(1) 证明在适当的初始条件下质点的运动轨道是一个圆；

(2) 如有心引力 $f(r) = -\frac{k}{r^{n}}$ ，其中 k 为正值常量，要此圆轨道运动是稳定的，对式中的 n 有何要求？
\end{problem}
\clearpage


\section{平方反比律的有心力作用下的运动}

\begin{problem}{4.2.1}
某个彗星围绕太阳做椭圆轨道运动，在远日点的速率为 $10 \, km/s$ ，在近日点的速率为 $80 \, km/s$ ，地球围绕太阳做圆轨道运动，速率为 $30 \, km/s$ ，半径为 $1.5 \times 10^{8} \, km$ ，求此彗星在远日点离太阳的距离。
\end{problem}
\clearpage


\begin{problem}{4.2.2}
假定在地球表面附近运行的人造地球卫星每 90min 转一周,月球表面附近运行的人造卫星也是每 90min 转一周,关于月球的组成你有何结论?
\end{problem}
\clearpage


\begin{problem}{4.2.3}
知道地球表面的重力加速度为 $9.8 \, m/s^{2}$ ，围绕地球的大圆周长为 $4 \times 10^{7} \, m$ ，又知道月球和地球的直径和质量之比分别为

\[
\frac {D _ {\mathrm{m}}}{D _ {\mathrm{e}}} = 0. 2 7, \quad \frac {M _ {\mathrm{m}}}{M _ {\mathrm{e}}} = 0. 0 1 2 3
\]

问从月球表面出发,逃离月球引力场所必须具有的相对于月球的最小速度多大(把月球视为惯性系)?
\end{problem}
\clearpage


\begin{problem}{4.2.4}
一长程火箭由半径为 R 的地球表面以速度 $v = v_{r} e_{r} + v_{\varphi} e_{\varphi}$ 发射，不计空气摩擦和地球转动，但需考虑精确的引力场，求出确定轨道所达到的最大高度 H 的方程（如图 4.8 所示），作最低级近似求解此方程，并证明对于 v 是垂直向上的情形，给出人们熟知的结果。
\begin{center}
\begin{tikzpicture}[>=Stealth, thick]
  % Earth circle of radius R
  \draw (0, 0) circle (1.5);
  \fill (0, 0) circle (1.4pt);

  % Radius arrow R
  \draw[->] (0, 0) -- (12:1.5) node[midway, below] {$R$};

  % Radial line to maximum height H (solid inside Earth, dashed above surface)
  \draw (0, 0) -- (62:1.5);
  \draw[dashed] (62:1.5) -- (62:2.6) node[pos=0.35, right=2pt] {$H$};

  % Rocket trajectory from surface up to smooth peak at (62:2.6) and curving down
  \coordinate (P) at (62:2.6);
  \draw (90:1.5) .. controls (76:2.15) and ($(P)+(152:0.28)$) .. (P)
        .. controls ($(P)+(-28:0.28)$) and (48:2.25) .. (40:1.82);
\end{tikzpicture}
\end{center}
\end{problem}
\clearpage


\begin{problem}{4.2.5}
设一宇宙飞船沿一个围绕太阳的椭圆轨道从地球发往火星, 近日点在地球绕太阳的轨道上, 远日点在火星绕太阳的轨道上, 如图 4.9 所示. 不考虑地球和火星对飞船的引力作用, 设地球和火星都绕太阳做圆周运动, 轨道半径分别为 $R_{1}, R_{2} = 1.5R_{1}$ . 求:

(1) 宇宙飞船的轨道方程；

(2) 先证明开普勒第三定律, 然后用它计算按上述轨道从地球到达火星所用的时间;

(3) 为了最节省燃料,从地球上应向什么方向发射?
\begin{center}
\begin{tikzpicture}[>=Stealth, thick]
  % Sun at origin
  \fill (0, 0) circle (1.5pt);
  \node[below=2pt] at (0, 0) {太阳};

  % Circular orbits of Earth (R1) and Mars (R2)
  \draw (0, 0) circle (1.35);
  \draw (0, 0) circle (2.15);

  % Spacecraft transfer ellipse tangent at (+1.35, 0) and (-2.15, 0)
  \draw (-0.40, 0) ellipse (1.75 and 0.78);

  % Radius vectors R1 and R2
  \draw[->] (0, 0) -- (52:1.35) node[pos=0.38, right=2pt] {$R_1$};
  \draw[->] (0, 0) -- (105:2.15) node[pos=0.78, left=2pt] {$R_2$};

  % Orbit labels and leader arrows
  \node[left] (earth) at (-2.25, 1.85) {地球轨道};
  \draw[->] (earth.east) -- (138:1.35);

  \node[left] at (-2.15, 1.05) {火星轨道};

  \node[left] (ship) at (-2.25, -1.55) {飞船轨道};
  \draw[->] (ship.east) -- (-1.35, -0.66);
\end{tikzpicture}
\end{center}
\end{problem}
\clearpage


\begin{problem}{4.2.6}
求上题所述的宇宙飞船需以多大的速度离开地球？仍不考虑地球的引力作用，也不考虑地球的自转。
\end{problem}
\clearpage


\begin{problem}{4.2.7}
若 4.2.5 题所述飞船飞离地球需考虑地球的引力作用,仍不考虑地球自转,求飞船从地球表面起航时需具有的速度.
\end{problem}
\clearpage


\begin{problem}{4.2.8}
质量为 $1.6 \times 10^{3} \mathrm{~kg}$ 的陨石在地球表面上方 $4.2 \times 10^{6} \mathrm{~m}$ 的圆形轨道上绕地球运动。它突然与另一轻得多的陨石发生正碰使动能损失 $2.0\%$ 。

(1) 什么物理定律适用于碰撞后的大陨石的运动?

(2) 描述碰撞后陨石轨道的形状；

(3) 求出陨石在碰撞后最接近地球的距离.
\end{problem}
\clearpage


\begin{problem}{4.2.9}
我国于 1971 年 3 月 3 日发射的科学实验卫星的重量为 221.0kg，近地点高度为 266km，远地点高度为 1826km，地球半径取为 6400km，求该卫星运行的周期。
\end{problem}
\clearpage


\begin{problem}{4.2.10}
一质量为 $m$ 的空间站沿半径为 $R$ 的圆轨道绕月球运动。为使空间站能在月球上登陆，沿轨道曲线方向向前发射一质量为 $m_{1}$ 的物体，使空间站在发射物体后，转过转角 $180^{\circ}$ 后刚好抵达月球表面。已知月球的质量为 $M_{\mathrm{m}}$ ，半径为 $R_{\mathrm{m}}$ 。求：

(1) 质量为 $m_{1}$ 的物体的发射速度；

(2) 从发射到空间站抵达月球所经历的时间.
\end{problem}
\clearpage


\begin{problem}{4.2.11}
假定你去访问一无大气、半径为 $R_0$ 的小行星，发现物体以速率 $v_0$ 水平抛出后，恰好能环绕该星体的表面做圆周运动。周 $v_0, R_0$ 表示下述各量：

(1) 物体逃离该小行星的速度；

(2) 从行星表面竖直上抛使最高点达到 $R_{0}$ 高度所需的上抛速率, 并求该物体达到

$\frac{1}{2} R_0$ 高度时的速率；

(3) 质量为 m 的物体在离行星表面为 y 时的势能；

(4) 使物体从行星表面升到 y 高度所需的上抛速率.
\end{problem}
\clearpage


\begin{problem}{4.2.12}
(1) 证明在按 $\frac{1}{r}$ 规律变化的有心引力势场中具有相同角动量的圆形轨道的半径等于抛物线轨道近心点离力心距离的两倍；

(2) 证明在(1)问中的有心力场中具有相同的角动量的圆形轨道和抛物线轨道相交处的速率, 后者是前者的 $\sqrt{2}$ 倍.
\end{problem}
\clearpage


\begin{problem}{4.2.13}
宇宙飞船绕一行星沿圆轨道飞行, 轨道半径为 R, 飞行速率为 $v_{0}$ , 要把轨道改为经过 B 点的椭圆轨道, 如图 4.10 所示, B 点位于以 O 为圆心、半径为 3R 的圆周上.

(1) 写出该椭圆的方程；

(2) 飞船在 A 点进入上述轨道时, 它的速度应增加多少?

(3) 从 A 点到 B 点的航程要用多少时间?

(4) 当飞船的位矢垂直于 AB 时, 速度的径向和横向分量多大?
\begin{center}
\begin{tikzpicture}[>=Stealth, thick]
  % Inner circular orbit of radius R centered at O = (0, 0)
  \draw (0, 0) circle (1.05);
  \node[right=1pt] at (0, 0) {$O$};
  \node[right=2pt] at (1.05, 0) {$A$};

  % Outer circular arc of radius 3R on the left side
  \draw (144:3.15) arc (144:214:3.15);
  \node[left=2pt] at (-3.15, 0) {$B$};

  % Transfer ellipse: periapsis at (+R, 0), apoapsis at (-3R, 0)
  % Center (-R, 0) = (-1.05, 0), semi-major a = 2R = 2.10, semi-minor b = sqrt(3)*R = 1.819
  \draw (-1.05, 0) ellipse (2.10 and 1.819);

  % Radius segments R and 3R
  \draw (0, 0) -- (135:1.05) node[midway, above right=0pt] {$R$};
  \draw (0, 0) -- (210:3.15) node[pos=0.45, above left=0pt] {$3R$};
\end{tikzpicture}
\end{center}
\end{problem}
\clearpage


\begin{problem}{4.2.14}
有两颗地球卫星 $M_{1}$ 及 $M_{2}$ 沿同一椭圆运动, 地球中心在此椭圆的一个焦点, 又设 $M_{1}M_{2}$ 相距不远, 可将圆弧 $\widehat{M_{1}M_{2}}$ 看作直线, 已知直线 $M_{1}M_{2}$ 的中点经近地点时, $M_{1}M_{2}=b$ , 近地点到地心的距离为 $R_{1}$ , 远地点到地心的距离为 $R_{2}$ . 求直线 $M_{1}M_{2}$ 的中点经远地点时两颗卫星的距离. 若取卫星处于近地点时 $\varphi=0$ , 则在任意 $\varphi$ 角时, 两颗卫星的距离多大?
\end{problem}
\clearpage


\begin{problem}{4.2.15}
一质点做椭圆轨道运动,偏心率为 e,力心在椭圆轨道的一焦点上.当质点运动到近心点时,力心突然移到椭圆的另一焦点上.求证以后椭圆轨道的偏心率为 $\frac{e(3+e)}{(1-e)}$ ,说明原来 e 值处在什么范围时,新的轨道将仍是椭圆、变成抛物线和变成双曲线
\end{problem}
\clearpage


\begin{problem}{4.2.16}
一行星绕质量为 $M$ 的恒星做圆周运动。恒星突然爆炸，以一个比行星的轨道速度大得多的速率喷射它的外层表壳，因此其质量损失可认为是瞬时的，恒星的剩余质量为 $M'$ ，仍比行星的质量 $m$ 大得多。求爆炸后行星轨道的偏心率忽略膨胀的壳体施加在行星上的力，可用偏心率 $e$ 与能量 $E$ 、角动量 $J$ 的下列关系：

\[
e ^ {2} = 1 + \frac {2 J ^ {2} E}{m ^ {3} \alpha^ {2}}
\]

恒星对行星的引力为 $F = -\frac{m\alpha}{r^2}$
\end{problem}
\clearpage


\begin{problem}{4.2.17}
已知一个由具有电荷 $e$ 的核和一个在半径为 $r_0$ 的圆形轨道上的单电子组成的氘原子的经典模型。突然这个核发射出一个电子，其电荷变为 $2e$ ，这个核发射的电子很快地逃逸了，我们可不考虑它，在轨道上的那个电子即有一个新状态。

(1) 求出发射后与发射前电子能量之比(规定无穷远处势能为零);

(2) 定性地描述新轨道；

(3) 求出新轨道的离核最近和最远的距离, 以 $r_{0}$ 为单位;

(4) 以 $r_{0}$ 表示出新轨道的长轴和短轴.
\end{problem}
\clearpage


\begin{problem}{4.2.18}
(1) 某彗星的轨道为抛物线, 其近日点为地球轨道(假定为圆形)半径的 $\frac{1}{n}$ , 则此彗星运行时, 在地球轨道内停留的时间为一年的 $\frac{2}{3\pi} \frac{n+2}{n} \sqrt{\frac{n-1}{2n}}$ 倍, 试证明之;

(2) 再证任何抛物线轨道的彗星停留在地球轨道(仍假定为圆形)内的最长时间为 $\frac{2}{3\pi}$ 年或约为76日.
\end{problem}
\clearpage


\begin{problem}{4.2.19}
一火箭自地球表面发射, 发射方向与铅直方向成 $\alpha$ 角, 火箭将按椭圆轨道运动, 若椭圆之长轴是地球半径的 $n(n < 2)$ 倍, 地球半径为 $R$ , 求火箭在空中飞行的时间.
\end{problem}
\clearpage


\begin{problem}{4.2.20}
证明做椭圆轨道运动的最大与最小速率的乘积为 $\left(\frac{2\pi a}{T}\right)^2$ ，其中 $T$ 是周期， $a$ 是长半轴。
\end{problem}
\clearpage


\section{有心力场中的散射}

\begin{problem}{4.3.1}
质量为 $m$ 的粒子以初速 $\bm{v}$ 从无穷远向其势能大小为 $\frac{\alpha}{r^n} (\alpha$ 为正值常量)的引力场运动, 求粒子被俘获的总截面.
\end{problem}
\clearpage


\begin{problem}{4.3.2}
一质量为 $m$ 的粒子以速度 $v_{0}$ 射向一固定散射中心。该散射中心施以排斥力 $F = \frac{1}{2} mv_{1}^{2}\delta (r - a)e_{r}$ ，其中 $\bm{e}_{r}$ 是由力心出发沿其径向的单位矢量， $a$ 是一固定半径，在该处有力作用， $v_{1}$ 是具有速度量纲的常数。

(1) 求势能；

(2) 证明若 $v_{0} < v_{1}$ , 则粒子不能进入半径 $r = a$ 的球内, 而被反射回去, 反射角等于入射角;

(3) 对 $v_{0}>v_{1}$ 、瞄准距离 $\rho=\frac{1}{2}a$ 的情形，画出粒子轨道的草图。
\end{problem}
\clearpage


\begin{problem}{4.3.3}
一质量为 $m$ 的粒子以能量 $E_0$ 、对 $O$ 点的角动量为 $J$ 被中心在 $O$ 点的势能 $V = -G(r)$ 的引力场所散射。

(1) 写出关于 $\frac{\mathrm{d}\alpha}{\mathrm{d}r}$ 的微分方程；

(2) 写出最接近距离 $r_{min}$ 所满足的方程.
\begin{center}
\begin{tikzpicture}[>=Stealth, line width=0.8pt]
  \def\R{1.6}
  \coordinate (O) at (0,0);
  \coordinate (Pin) at (150:\R);
  \coordinate (Pout) at (60:\R);

  % Main circle
  \draw (O) circle (\R);

  % Radial and vertical dashed lines
  \draw[dashed] (O) -- (0, 2.35);
  \draw[dashed] (O) -- (150:2.75);
  \draw[dashed] (O) -- (60:2.55);

  % Horizontal dashed line from Pin to dimension ticks
  \draw[dashed] (Pin) -- (0.5, {0.5*\R});
  \draw (0.15, 0) -- (0.6, 0);
  \draw (0.3, {0.5*\R}) -- (0.6, {0.5*\R});
  \draw[<->] (0.45, 0) -- (0.45, {0.5*\R}) node[midway, right, inner sep=2pt] {$\frac{1}{2}a$};

  % Incoming ray v_0
  \draw[->] (-2.7, {0.5*\R}) -- (Pin);
  \node[below] at (-2.05, {0.5*\R}) {$v_0$};

  % Angle 30 degrees at Pin
  \draw (Pin) ++(180:0.65) arc (180:150:0.65);
  \node at ($(Pin)+(165:0.98)$) {$30^\circ$};

  % Chord Pin -> Pout with mid-arrow
  \draw[postaction={decorate}, decoration={markings, mark=at position 0.55 with {\arrow{>}}}]
    (Pin) -- (Pout);

  % Outgoing ray v''
  \draw[->] (Pout) -- ++(40:1.15) node[above right, inner sep=1pt] {$v''$};

  % Angle alpha between outgoing normal and v''
  \draw[->] ($(Pout)+(60:1.1)+(135:0.35)$) node[above left, inner sep=1pt] {$\alpha$}
    to[bend left=25] ($(Pout)+(46:0.75)$);

  % Angle theta between radial dashed line (to O) and chord
  \draw (Pin) ++(-30:0.65) arc (-30:15:0.65);
  \node at ($(Pin)+(-10:0.9)$) {$\theta$};

  % Velocity label v' with curved pointer arrow from chord
  \draw[->] ($(Pin)!0.22!(O)$) to[out=70, in=-90] (-0.9, 1.95) node[above] {$v'$};
\end{tikzpicture}

\end{center}
\begin{center}
\begin{tikzpicture}[>=Stealth, line width=0.8pt]
  \coordinate (O) at (0,0);
  \coordinate (L) at (-3.2, 0);

  % Horizontal reference line and center O
  \draw (L) -- (O) node[below right, inner sep=2pt] {$O$};

  % Scattering trajectory curve
  \draw (-2.7, 1.85) .. controls (-0.6, 1.45) and (1.15, 0.95) .. (1.28, 0.05)
    .. controls (1.35, -0.55) and (1.1, -1.15) .. (0.8, -1.75);

  % Position vector r at angle alpha
  \coordinate (P) at (-1.38, 1.54);
  \draw (O) -- (P) node[midway, above right, inner sep=1pt] {$r$};

  % Angle alpha arc
  \draw (180:0.65) arc (180:132:0.65);
  \node at (156:0.9) {$\alpha$};
\end{tikzpicture}

\end{center}
\end{problem}
\clearpage


\begin{problem}{4.3.4}
一束很宽的平行小粒子束, 粒子质量为 m, 不带电, 从空间射向月球, 相对于月球的初速为 $v_{0}$ , 忽略地球和太阳的存在, 求用月球的半径 R、从月球表面的逃逸速度 $v_{esc}$ 和 $v_{0}$ 表示的碰击月球的碰撞截面 $\sigma$ .
\end{problem}
\clearpage


\begin{problem}{4.3.5}
一彗星以初速 $v_{0}$ 向太阳运动, 太阳的质量为 $M$ 、半径为 $R$ , 可视为静止的, 并略去所有其他的星体. 求彗星和太阳相撞的总截面.
\end{problem}
\clearpage


\begin{problem}{4.3.6}
一个原子与一个离子的相互作用可由势能 $V(r) = -cr^{-4}$ 给出，其中 $c = \frac{e^2p_a^2}{2}$ ， $e$ 是离子电荷， $p_a$ 是原子的极化率。

(1) 作出有效势能 $U(r)$ 的草图；

(2) 求出用角动量 J 表达的有效势能的极大值 $U_{0}$ ;

(3) 设离子远比原子轻, 求离子击中原子(即深入到 $r = 0$ ) 的横截面 $\sigma(v_0)$ .
\end{problem}
\clearpage


\begin{problem}{4.3.7}
求一个粒子被半径为 $a$ 的刚球弹性散射的微分散射截面和总散射截面。
\end{problem}
\clearpage


\begin{problem}{4.3.8}
一质量为 $m$ 、半径为 $r$ 的小弹子球被一质量为 $M(m \ll M)$ 、半径为 $R$ 的大弹子球弹性散射，求微分散射截面和总散射截面。
\end{problem}
\clearpage


\begin{problem}{4.3.9}
质量为 $m$ 的粒子以速度 $v_{0}$ 从无穷远发射被大小为 $\frac{k}{r^2} (k$ 为常量)的固定的斥力中心所散射。如粒子不受斥力作用，粒子将以与斥力中心距离为 $b$ 的轨迹通过，求：

(1) 粒子接近斥力中心的最近距离；

(2) 实际产生的偏转角；

(3) 对于一均匀的粒子束受到该势场散射的微分散射截面 $\frac{d\sigma}{d\Omega}$ .
\end{problem}
\clearpage


\begin{problem}{4.3.10}
(1) 质量为 $m$ 的粒子在势场 $V(r) = \frac{k}{r^2} (k > 0$ , 为常量) 中运动, 考虑在 $xy$ 平面内运动, 用极坐标 $r, \varphi$ , 求出 $r$ 作为 $\varphi$ 、角动量 $J$ 和能量 $E$ 的函数关系;

(2) 用(1)问中的结果讨论在这个势场中的散射, 求微分散射截面 $\sigma(\theta, E)$ , 其中 $\theta$ 为散射角.
\end{problem}
\clearpage


\begin{problem}{4.3.11}
具有动能为 $T$ (非相对论) 的带电 $\pi$ 介子 ( $\pi^{+}$ 或 $\pi^{-}$ ) 和电荷为 $Ze$ 、有效半径为 $b$ 的重核。考虑经典力学的图像，当 $\pi$ 介子距核的最近距离等于或小于 $b$ 时， $\pi$ 介子就会打到核上。忽略核的反冲，证明 $\pi$ 介子的碰撞截面为

\[
\sigma = \pi b ^ {2} \Big (\frac {T - V}{T} \Big) \qquad \text { 对于 } \pi^ {+}
\]

\[
\sigma = \pi b ^ {2} \Big (\frac {T + V}{T} \Big) \qquad \text {对于} \pi^ {-}
\]

其中

\[
V = \frac {Z e ^ {2}}{4 \pi \varepsilon_ {0} b}.
\]
\end{problem}
\clearpage


\begin{problem}{4.3.12}
在核物理中常遇到下列球形势阱：

\[
V (r) = \left\{ \begin{array}{l l} - V _ {0} & r \leqslant a \\ 0 & r > a \end{array} \right.
\]

(1) 证明: 在经典力学中被上述势产生的散射如同光线被相对折射率 $n = \sqrt{\frac{E + V_0}{E}}$ 、半径为 $a$ 的球的折射 (这等效说明为什么既能用惠更斯原理的波动说又可用牛顿力学的微粒说解释折射现象);

(2) 证明微分散射截面为

\[
\sigma (E, \theta) = \left\{ \begin{array}{l l} \frac {n ^ {2} a ^ {2}}{4 \cos \frac {\theta}{2}} \frac {(n \cos \frac {\theta}{2} - 1) (n - \cos \frac {\theta}{2})}{(1 + n ^ {2} - 2 n \cos \frac {\theta}{2}) ^ {2}} & 0 \leqslant \theta \leqslant 2 \arccos \frac {1}{n} \\ 0 & \theta \geqslant 2 \arccos \frac {1}{n} \end{array} \right.
\]

(3) 求总散射截面, 可能用到下列积分公式:

\[
\begin{array}{r l} & {\int \frac {x \mathrm{d} x}{(a + b x) ^ {2}} = \frac {1}{b ^ {2}} \Big [ \ln (a + b x) + \frac {a}{a + b x} \Big ]} \\ & {\int \frac {x ^ {2} \mathrm{d} x}{(a + b x) ^ {2}} = \frac {1}{b ^ {3}} \Big [ a + b x - 2 a \ln (a + b x) - \frac {a ^ {2}}{a + b x} \Big ]} \end{array}
\]
\end{problem}
\clearpage


\begin{problem}{4.3.13}
一束质量为 $m$ 、能量为 $E$ 的粒子垂直入射在厚度为 $d$ 、每单位体积内包含 $n$ 个固定散射中心的金属薄片上。一个粒子与一个散射中心间的力是大小为 $\frac{k}{r^2}$ ( $k$ 为常量) 的斥力。问束中粒子有多大部分被散射到 $\theta$ 大于 $\alpha$ 的立体角内？
\end{problem}
\clearpage


\begin{problem}{4.3.14}
一束质量为 $m$ 、能量为 $E$ 的粒子垂直入射在上题所述的金属薄片上，离粒子束通过薄片所在的 $O$ 点距离为 $r$ 的地方，放置一面积为 $A(A \ll r^2)$ 的粒子探测器，从 $O$ 点到探测器的线与粒子入射方向间夹角为 $\alpha$ (如图4.21所示), 若射到金属薄片上的粒子束流密度为 $I$ 、问单位时间内有多少粒子射中探测器?
\begin{center}
\begin{tikzpicture}[>=Stealth, line width=0.8pt]
  \coordinate (O) at (0,0);

  % Thin foil horizontal line
  \draw (-2.3, 0) -- (2.3, 0);
  \node[below left, inner sep=2pt] at (O) {$O$};

  % Incoming particle beam
  \draw[postaction={decorate}, decoration={markings, mark=at position 0.55 with {\arrow{>}}}]
    (0, 2.0) -- (O);
  \node[right, align=center] at (0.12, 1.15) {粒\\子\\束};

  % Vertical dashed continuation
  \draw[dashed] (O) -- (0, -2.0);

  % Scattered ray to detector at angle alpha
  \def\ang{-58}
  \coordinate (Dfront) at (\ang:2.1);
  \draw (O) -- (Dfront) node[midway, above right, inner sep=1pt] {$r$};

  % Angle alpha
  \draw (-90:0.6) arc (-90:\ang:0.6);
  \node at ({(-90+\ang)/2}:0.9) {$\alpha$};

  % Detector rectangle aligned with ray
  \draw[fill=white] (Dfront) -- ++(\ang+90:0.16) -- ++(\ang:0.65) -- ++(\ang-90:0.32) -- ++(\ang+180:0.65) -- cycle;
  \node[above right, inner sep=2pt] at ($(Dfront)+(\ang+90:0.16)$) {$A$};
  \node[below, inner sep=2pt] at ($(Dfront)+(\ang:0.78)+(-90:0.15)$) {探测器};
\end{tikzpicture}

\end{center}
\end{problem}
\clearpage


\begin{problem}{4.3.15}
粒子束遇到由许多散射中心组成的靶子．设 $I(0)$ 是入射时的粒子束强度，进入靶后穿透 $x$ 距离处的强度为 $I(x)$ ．设束中一个粒子被一个散射中心散射的总散射截面为 $\sigma$ ，证明：

(1) $I(x)=I(0)e^{-n\sigma x},n$ 为单位体积内的散射中心数；

(2) 粒子被散射前运行的平均距离为 $\frac{1}{n\sigma}$ .
\end{problem}
\clearpage

